\subsection{\texorpdfstring{第~\ref{sec:differential-forms}~节}{微分形式}}\label{subsec:solutions-pullback-and-exterior-derivative}

\begin{description}[font=\normalfont,leftmargin=2em,labelwidth=1.5em]
\RGsolution{ex:prove-pullback-of-differential-forms}{pullback-of-differential-forms}
性质 (i) 显然成立, 性质 (iii) 是 Proposition~\ref{prop:pullback-of-tensors} 的直接推论. 为证明 (ii), 注意给定 $p\in M$ 和 $v_1,\ldots,v_k\in T_pM$, 有
\begin{align*}
  (f^*(g\alpha))_p(v_1,\ldots,v_k)
  &=(g\alpha)_{f(p)}((\mathrm{d}f)_pv_1,\ldots,(\mathrm{d}f)_pv_k)\\
  &=g(f(p))\alpha_{f(p)}((\mathrm{d}f)_pv_1,\ldots,(\mathrm{d}f)_pv_k)\\
  &=(g\circ f)(p)(f^*\alpha)_p(v_1,\ldots,v_k)\\
  &=((f^*g)(f^*\alpha))_p(v_1,\ldots,v_k).
\end{align*}
最后, (iv) 由下式得出:
\begin{align*}
  (g^*f^*\alpha)_p(v_1,\ldots,v_k)
  &=(f^*\alpha)_{g(p)}((\mathrm{d}g)_pv_1,\ldots,(\mathrm{d}g)_pv_k)\\
  &=\alpha_{f(g(p))}((\mathrm{d}f)_{g(p)}(\mathrm{d}g)_pv_1,\ldots,
                          (\mathrm{d}f)_{g(p)}(\mathrm{d}g)_pv_k)\\
  &=\alpha_{(f\circ g)(p)}((\mathrm{d}(f\circ g))_pv_1,\ldots,(\mathrm{d}(f\circ g))_pv_k)\\
  &=((f\circ g)^*\alpha)_p(v_1,\ldots,v_k).
\end{align*}

\RGsolution{ex:properties-of-exterior-derivative-on-manifolds}{properties-of-exterior-derivative-on-manifolds}
设 $M$ 是光滑流形.
\begin{enumerate}[label=(\roman*),ref=\roman*,leftmargin=*]
\item $\mathrm{d}(\omega_1+\omega_2)$ 是对每个参数化 $\varphi_\alpha:U_\alpha\to M$ 都以 $\mathrm{d}(\omega_1+\omega_2)_\alpha$ 为其局部表示的形式. 此外, 由于
\[
  (\omega_1+\omega_2)_\alpha=\varphi_\alpha^*(\omega_1+\omega_2)
  =\varphi_\alpha^*\omega_1+\varphi_\alpha^*\omega_2=(\omega_1)_\alpha+(\omega_2)_\alpha,
\]
由 Proposition~\ref{prop:properties-of-exterior-derivative}(i) 得
\begin{align*}
  \mathrm{d}(\omega_1+\omega_2)_\alpha
  &=\mathrm{d}((\omega_1)_\alpha+(\omega_2)_\alpha)
  =\mathrm{d}(\omega_1)_\alpha+\mathrm{d}(\omega_2)_\alpha\\
  &=\varphi_\alpha^*\mathrm{d}\omega_1+\varphi_\alpha^*\mathrm{d}\omega_2
  =\varphi_\alpha^*(\mathrm{d}\omega_1+\mathrm{d}\omega_2),
\end{align*}
故 $\mathrm{d}(\omega_1+\omega_2)=\mathrm{d}\omega_1+\mathrm{d}\omega_2$.
\item 设 $\omega$ 是 $k$-形式. 则 $\mathrm{d}(\omega\wedge\gamma)$ 是对每个参数化 $\varphi_\alpha:U_\alpha\to M$ 都以 $\mathrm{d}(\omega\wedge\gamma)_\alpha$ 为其局部表示的形式. 此外, 由 Proposition~\ref{prop:pullback-of-differential-forms}(iii),
\[
  (\omega\wedge\gamma)_\alpha
  =\varphi_\alpha^*(\omega\wedge\gamma)
  =\varphi_\alpha^*\omega\wedge\varphi_\alpha^*\gamma=\omega_\alpha\wedge\gamma_\alpha.
\]
因此, 由 Proposition~\ref{prop:properties-of-exterior-derivative}(ii),
\begin{align*}
  \mathrm{d}(\omega\wedge\gamma)_\alpha
  &=\mathrm{d}(\omega_\alpha\wedge\gamma_\alpha)
  =\mathrm{d}\omega_\alpha\wedge\gamma_\alpha+(-1)^k\omega_\alpha\wedge\mathrm{d}\gamma_\alpha\\
  &=\varphi_\alpha^*(\mathrm{d}\omega)\wedge\varphi_\alpha^*\gamma
  +(-1)^k\varphi_\alpha^*\omega\wedge\varphi_\alpha^*(\mathrm{d}\gamma)\\
  &=\varphi_\alpha^*(\mathrm{d}\omega\wedge\gamma+(-1)^k\omega\wedge\mathrm{d}\gamma)\\
  &=(\mathrm{d}\omega\wedge\gamma+(-1)^k\omega\wedge\mathrm{d}\gamma)_\alpha,
\end{align*}
所求结论随即成立.
\item 对每个参数化 $\varphi_\alpha:U_\alpha\to M$, 形式 $\mathrm{d}(\mathrm{d}\omega)$ 的局部表示为 $\mathrm{d}(\mathrm{d}\omega)_\alpha=\mathrm{d}(\mathrm{d}\omega_\alpha)$. 由 Proposition~\ref{prop:properties-of-exterior-derivative}(iii), 有 $\mathrm{d}(\mathrm{d}\omega_\alpha)=0$, 因而 $\mathrm{d}(\mathrm{d}\omega)=0$.
\item 设 $f:M\to N$ 是两个流形间的光滑映射, $\omega$ 是 $N$ 上的形式. 对每个参数化 $\varphi_\alpha:U_\alpha\to M$, $\mathrm{d}(f^*\omega)$ 的局部表示是 $\mathrm{d}(f^*\omega)_\alpha$. 此外, 由 Proposition~\ref{prop:pullback-of-differential-forms}(iv),
\[
  (f^*\omega)_\alpha=\varphi_\alpha^*(f^*\omega)=(f\circ\varphi_\alpha)^*\omega,
\]
因而
\[
  \mathrm{d}(f^*\omega)_\alpha=\mathrm{d}((f\circ\varphi_\alpha)^*\omega).
\]
另一方面, 对每个参数化 $\varphi_\alpha:U_\alpha\to M$, 形式 $f^*(\mathrm{d}\omega)$ 的局部表示是
\[
  \varphi_\alpha^*(f^*(\mathrm{d}\omega))=(f\circ\varphi_\alpha)^*(\mathrm{d}\omega).
\]
若 $\psi_\beta:V\to N$ 是 $N$ 的一个参数化, 则在 $\psi_\beta(V)$ 上有
\[
  \mathrm{d}\omega=(\psi_\beta^{-1})^*\mathrm{d}\omega_\beta
  =(\psi_\beta^{-1})^*\mathrm{d}(\psi_\beta^*\omega),
\]
所以由 Proposition~\ref{prop:properties-of-exterior-derivative}, 在 $\varphi_\alpha^{-1}(f^{-1}(\psi_\beta(V)))$ 上有
\begin{align*}
  \varphi_\alpha^*(f^*(\mathrm{d}\omega))
  &=(f\circ\varphi_\alpha)^*(\psi_\beta^{-1})^*\mathrm{d}\omega_\beta
  =(\psi_\beta^{-1}\circ f\circ\varphi_\alpha)^*\mathrm{d}\omega_\beta\\
  &=\mathrm{d}((\psi_\beta^{-1}\circ f\circ\varphi_\alpha)^*\omega_\beta)
  =\mathrm{d}((f\circ\varphi_\alpha)^*(\psi_\beta^{-1})^*\omega_\beta)\\
  &=\mathrm{d}((f\circ\varphi_\alpha)^*\omega).
\end{align*}
由于 $\psi_\beta:V\to N$ 是任取的, 所求结论成立.
\end{enumerate}

\RGsolution{ex:homotopy-operator-and-poincare-lemma}{homotopy-operator-and-poincare-lemma}
\begin{enumerate}[label=(\alph*),ref=\alph*,leftmargin=*]
\item 有
\[
  \widetilde f^*\omega=\mathrm{d}t\wedge\sum_I(a_I\circ\widetilde f)\,\mathrm{d}f^I
   +\sum_J(b_J\circ\widetilde f)\,\mathrm{d}f^J,
\]
因而
\[
  Q(\widetilde f^*\omega)=\sum_I\left(\int_{t_0}^t(a_I\circ\widetilde f)\,\mathrm{d}s\right)\mathrm{d}f^I.
\]
另一方面, 由于 $\widetilde f$ 在第一个坐标上是恒等映射, 先施行 $\widetilde f$ 再施行被积函数的复合操作与对 $t$ 积分可交换. 因此,
\begin{align*}
  \widetilde f^*(Q(\omega))
  &=\sum_I\left(\left(\int_{t_0}^t a_I\,\mathrm{d}s\right)\circ\widetilde f\right)\mathrm{d}f^I\\
  &=\sum_I\left(\int_{t_0}^t(a_I\circ\widetilde f)\,\mathrm{d}s\right)\mathrm{d}f^I=Q(\widetilde f^*\omega).
\end{align*}
\item 每个参数化 $\varphi:U\to M$ 都通过 $\widetilde\varphi=\operatorname{id}\times\varphi$ 给出一个参数化 $\widetilde\varphi:\mathbb{R}\times U\to\mathbb{R}\times M$. 给定 $k$-形式 $\omega\in\Omega^k(\mathbb{R}\times M)$, 将 $Q(\omega)\in\Omega^{k-1}(\mathbb{R}\times M)$ 定义为在参数化 $\widetilde\varphi_\alpha$ 下局部表示为 $(Q(\omega))_\alpha=Q(\omega_\alpha)$ 的 $(k-1)$-形式. 即在 $\mathbb{R}\times\varphi_\alpha(U_\alpha)$ 上, 有 $Q(\omega)=(\widetilde\varphi_\alpha^{-1})^*(Q(\omega_\alpha))$. 为检验这个定义的一致性, 考虑另一个参数化 $\varphi_\beta:U_\beta\to M$, 使 $W:=\varphi_\alpha(U_\alpha)\cap\varphi_\beta(U_\beta)\ne\varnothing$. 设 $f:\varphi_\alpha^{-1}(W)\to\varphi_\beta^{-1}(W)$ 是由 $f=\varphi_\beta^{-1}\circ\varphi_\alpha$ 给出的微分同胚. 则在 $\mathbb{R}\times W$ 上,
\begin{align*}
  (\widetilde\varphi_\alpha^{-1})^*(Q(\omega_\alpha))
  &=(\widetilde\varphi_\alpha^{-1})^*(Q(\widetilde f^*\omega_\beta))
  =(\widetilde\varphi_\alpha^{-1})^*(\widetilde f^*(Q(\omega_\beta)))\\
  &=(\widetilde f\circ\widetilde\varphi_\alpha^{-1})^*(Q(\omega_\beta))
  =(\widetilde\varphi_\beta^{-1})^*(Q(\omega_\beta)).
\end{align*}
为了验证在给定点 $(t,p)\in\mathbb{R}\times M$ 处有 $\widetilde f^*\circ Q=Q\circ\widetilde f^*$, 考虑 $p$ 和 $f(p)$ 附近的参数化 $\varphi:U\to M$ 和 $\psi:V\to N$, 设 $\widehat f=\psi^{-1}\circ f\circ\varphi$ 为 $f$ 相应的局部表示. 则在 $p$ 处有
\begin{align*}
  \widetilde f^*(Q(\omega))
  &=\widetilde f^*(\widetilde\psi^{-1})^*(Q(\widetilde\psi^*\omega))
  =(\widetilde\psi^{-1}\circ\widetilde f)^*(Q(\widetilde\psi^*\omega))\\
  &=(\widetilde{\widehat f}\circ\widetilde\varphi^{-1})^*(Q(\widetilde\psi^*\omega))
  =(\widetilde\varphi^{-1})^*\widetilde{\widehat f}^{\,*}(Q(\widetilde\psi^*\omega))\\
  &=(\widetilde\varphi^{-1})^*(Q(\widetilde{\widehat f}^{\,*}(\widetilde\psi^*\omega)))
  =(\widetilde\varphi^{-1})^*(Q((\widetilde\psi\circ\widetilde{\widehat f})^*\omega))\\
  &=(\widetilde\varphi^{-1})^*(Q((\widetilde f\circ\widetilde\varphi)^*\omega))
  =(\widetilde\varphi^{-1})^*(Q(\widetilde\varphi^*\widetilde f^*\omega))\\
  &=Q(\widetilde f^*\omega).
\end{align*}
类似地, $M$ 上 $Q$ 的线性性由 $\mathbb{R}^n$ 上 $Q$ 的线性性得出.
\item 注意, 在局部坐标下 $Q$ 的公式化为 $\mathbb{R}^n$ 中的公式. 若
\[
  \omega=\mathrm{d}t\wedge\sum_I a_I(t,x)\,\mathrm{d}x^I+\sum_J b_J(t,x)\,\mathrm{d}x^J,
\]
则
\[
  i_{t_0}^*\omega=\sum_J b_J(t_0,x)\,\mathrm{d}x^J
\]
(因为 $t_0$ 为常数, 所以 $\mathrm{d}t_0=0$), 而 $\pi^*i_{t_0}^*\omega$ 由同一表达式给出. 因此
\begin{align*}
  \omega-\pi^*i_{t_0}^*\omega
  &=\mathrm{d}t\wedge\sum_I a_I(t,x)\,\mathrm{d}x^I
    +\sum_J(b_J(t,x)-b_J(t_0,x))\,\mathrm{d}x^J.
\end{align*}
另一方面,
\[
  \mathrm{d}(Q(\omega))=\mathrm{d}t\wedge\sum_I a_I(t,x)\,\mathrm{d}x^I
    +\sum_I\sum_i\left(\int_{t_0}^t\frac{\partial a_I}{\partial x^i}\,\mathrm{d}s\right)\mathrm{d}x^i\wedge\mathrm{d}x^I
\]
且
\begin{align*}
  \mathrm{d}\omega
  &=-\mathrm{d}t\wedge\sum_I\sum_i\frac{\partial a_I}{\partial x^i}\mathrm{d}x^i\wedge\mathrm{d}x^I\\
  &\quad+\mathrm{d}t\wedge\sum_J\frac{\partial b_J}{\partial t}\mathrm{d}x^J
       +\sum_J\sum_i\frac{\partial b_J}{\partial x^i}\mathrm{d}x^i\wedge\mathrm{d}x^J,
\end{align*}
由此得到
\begin{align*}
  Q(\mathrm{d}\omega)
  &=-\sum_I\sum_i\left(\int_{t_0}^t\frac{\partial a_I}{\partial x^i}\,\mathrm{d}s\right)\mathrm{d}x^i\wedge\mathrm{d}x^I\\
  &\quad+\sum_J\left(\int_{t_0}^t\frac{\partial b_J}{\partial s}\,\mathrm{d}s\right)\mathrm{d}x^J.
\end{align*}
从而,
\begin{align*}
  \mathrm{d}(Q(\omega))+Q(\mathrm{d}\omega)
  &=\mathrm{d}t\wedge\sum_I a_I(t,x)\,\mathrm{d}x^I
    +\sum_J(b(t,x)-b(t_0,x))\,\mathrm{d}x^J\\
  &=\omega-\pi^*i_{t_0}^*\omega.
\end{align*}
\item\label{sol:homotopy-operator-cylinder-cohomology} 由于 $\pi\circ i_{t_0}=\operatorname{id}_M$, 有 $i_{t_0}^*\circ\pi^*=\operatorname{id}$, 因而 $i_{t_0}^\sharp\circ\pi^\sharp=\operatorname{id}$. 另一方面, 若 $\omega$ 是闭形式, 则
\[
  \omega-\pi^*i_{t_0}^*\omega=\mathrm{d}(Q(\omega)),
\]
这意味着 $\omega$ 和 $\pi^*i_{t_0}^*\omega$ 在同一个上同调类中. 因此 $\pi^\sharp\circ i_{t_0}^\sharp=\operatorname{id}$.
\item 由 (\ref{sol:homotopy-operator-cylinder-cohomology}), 有 $H^k(\mathbb{R}^n)=H^k(\mathbb{R}^{n-1})=\cdots=H^k(\mathbb{R})$. 当 $k>1$ 时, 已知 $H^k(\mathbb{R})=0$. 另一方面, $\mathbb{R}$ 上的任意 1-形式必是恰当的, 因为
\[
  a(t)\,\mathrm{d}t=\mathrm{d}\left(\int_0^t a(s)\,\mathrm{d}s\right).
\]
因此 $H^1(\mathbb{R})=0$. 所以对于所有 $k>0$, 都有 $H^k(\mathbb{R}^n)=0$.
\item 有 $f=H\circ i_{t_0}$ 和 $g=H\circ i_{t_1}$, 所以 $f^\sharp=i_{t_0}^\sharp\circ H^\sharp$ 且 $g^\sharp=i_{t_1}^\sharp\circ H^\sharp$. 因为 $i_{t_0}^\sharp$ 和 $i_{t_1}^\sharp$ 都是 $\pi^\sharp$ 的逆映射, 所以 $i_{t_0}^\sharp=i_{t_1}^\sharp$. 从而 $f^\sharp=g^\sharp$.
\item 由 $H(t,x)=(1-t)x$ 给出的映射 $H:\mathbb{R}\times\mathbb{R}^n\to\mathbb{R}^n$ 显然是恒等映射 $f(x)=H(0,x)=x$ 与常值映射 $g(x)=H(1,x)=0$ 之间的同伦.
\item 由于 $M$ 可缩, 恒等映射 $\operatorname{id}:M\to M$ 光滑同伦于常值映射 $g:M\to M$, 所以 $\operatorname{id}^\sharp=g^\sharp$. 利用局部坐标, 立即可验证, 若 $\omega\in\Omega^k(M)$ 是 $k>0$ 的 $k$-形式, 则 $g^*\omega=0$, 因而 $g^\sharp:H^k(M)\to H^k(M)$ 为零映射. 但 $g^\sharp=\operatorname{id}^\sharp$ 是 $H^k(M)$ 上的恒等映射, 所以必有 $H^k(M)=0$.
\end{enumerate}

\RGsolution{ex:lie-derivative-of-differential-form}{lie-derivative-of-differential-form}
\begin{enumerate}[label=(\alph*),ref=\alph*,leftmargin=*]
\item 有
\begin{align*}
  L_X(\omega_1\wedge\omega_2)
  &=\left.\frac{\mathrm{d}}{\mathrm{d}t}(\psi_t^*(\omega_1\wedge\omega_2))\right|_{t=0}
   =\left.\frac{\mathrm{d}}{\mathrm{d}t}((\psi_t^*\omega_1)\wedge(\psi_t^*\omega_2))\right|_{t=0}\\
  &=\left[\left.\frac{\mathrm{d}}{\mathrm{d}t}(\psi_t^*\omega_1)\right|_{t=0}\right]\wedge\omega_2
   +\omega_1\wedge\left[\left.\frac{\mathrm{d}}{\mathrm{d}t}(\psi_t^*\omega_2)\right|_{t=0}\right]\\
  &=(L_X\omega_1)\wedge\omega_2+\omega_1\wedge(L_X\omega_2)
\end{align*}
(这里使用了楔积的双线性).
\item 类似地,
\begin{align*}
  L_X(\mathrm{d}\omega)
  &=\left.\frac{\mathrm{d}}{\mathrm{d}t}(\psi_t^*(\mathrm{d}\omega))\right|_{t=0}
   =\left.\frac{\mathrm{d}}{\mathrm{d}t}(\mathrm{d}(\psi_t^*\omega))\right|_{t=0}\\
  &=\mathrm{d}\left[\left.\frac{\mathrm{d}}{\mathrm{d}t}(\psi_t^*\omega)\right|_{t=0}\right]
   =\mathrm{d}(L_X\omega)
\end{align*}
(这里使用了外微分的线性性).
\item 用归纳法证明此公式. 先考虑 $\omega=f\,\mathrm{d}g$ 的情形, 其中 $f,g\in\Omega^0(M)$ 是光滑函数. 这时, 由上面证明的性质得
\[
  L_X\omega=(L_X f)\,\mathrm{d}g+f\,\mathrm{d}(L_Xg)
  =(X\cdot f)\,\mathrm{d}g+f\,\mathrm{d}(X\cdot g).
\]
另一方面,
\begin{align*}
  \iota(X)\mathrm{d}\omega+\mathrm{d}(\iota(X)\omega)
  &=\iota(X)(\mathrm{d}f\wedge\mathrm{d}g)+\mathrm{d}(f(X\cdot g))\\
  &=\iota(X)(\mathrm{d}f\otimes\mathrm{d}g-\mathrm{d}g\otimes\mathrm{d}f)
    +(X\cdot g)\,\mathrm{d}f+f\,\mathrm{d}(X\cdot g)\\
  &=(X\cdot f)\,\mathrm{d}g-(X\cdot g)\,\mathrm{d}f+(X\cdot g)\,\mathrm{d}f+f\,\mathrm{d}(X\cdot g)\\
  &=L_X\omega.
\end{align*}
接下来证明, 若 Cartan 公式对 $\omega$ 和 $\eta$ 成立, 则它也对 $\omega\wedge\eta$ 成立. 由于任何形式在局部都可由形如 $f\,\mathrm{d}g$ 的 1-形式的楔积得到, 这就完成了证明. 若 $\omega$ 是 $k$-形式, 则
\begin{align*}
  &\iota(X)\mathrm{d}(\omega\wedge\eta)+\mathrm{d}(\iota(X)(\omega\wedge\eta))\\
  &\quad=\iota(X)(\mathrm{d}\omega\wedge\eta+(-1)^k\omega\wedge\mathrm{d}\eta)
     +\mathrm{d}((\iota(X)\omega)\wedge\eta+(-1)^k\omega\wedge(\iota(X)\eta))\\
  &\quad=(\iota(X)\mathrm{d}\omega)\wedge\eta
     +(-1)^{k+1}\mathrm{d}\omega\wedge(\iota(X)\eta)
     +(-1)^k(\iota(X)\omega)\wedge\mathrm{d}\eta\\
  &\qquad+\omega\wedge(\iota(X)\mathrm{d}\eta)
     +\mathrm{d}(\iota(X)\omega)\wedge\eta
     +(-1)^{k-1}(\iota(X)\omega)\wedge\mathrm{d}\eta\\
  &\qquad+(-1)^k\mathrm{d}\omega\wedge(\iota(X)\eta)+\omega\wedge\mathrm{d}(\iota(X)\eta)\\
  &\quad=(\iota(X)\mathrm{d}\omega)\wedge\eta+\omega\wedge(\iota(X)\mathrm{d}\eta)
        +\mathrm{d}(\iota(X)\omega)\wedge\eta+\omega\wedge\mathrm{d}(\iota(X)\eta)\\
  &\quad=(L_X\omega)\wedge\eta+\omega\wedge(L_X\eta)=L_X(\omega\wedge\eta)
\end{align*}
(这里使用了 Exercise~\ref{ex:tensors}(\ref{ex:contraction-of-alternating-tensors})).
\item 首先注意,
\begin{align*}
  &\bigl(\psi_t^*(\iota(Y)\omega)\bigr)_p(v_1,\ldots,v_k)
   =(\iota(Y)\omega)_{\psi_t(p)}((\mathrm{d}\psi_t)_pv_1,\ldots,(\mathrm{d}\psi_t)_pv_k)\\
  &\quad=\omega_{\psi_t(p)}(Y_{\psi_t(p)},(\mathrm{d}\psi_t)_pv_1,\ldots,(\mathrm{d}\psi_t)_pv_k)\\
  &\quad=\omega_{\psi_t(p)}((\mathrm{d}\psi_t)_p(\mathrm{d}\psi_{-t})_{\psi_t(p)}Y_{\psi_t(p)},
                            (\mathrm{d}\psi_t)_pv_1,\ldots,(\mathrm{d}\psi_t)_pv_k)\\
  &\quad=(\psi_t^*\omega)_p((\mathrm{d}\psi_{-t})_{\psi_t(p)}Y_{\psi_t(p)},v_1,\ldots,v_k),
\end{align*}
即
\[
  \psi_t^*(\iota(Y)\omega)=\iota((\psi_{-t})_*Y)(\psi_t^*\omega).
\]
在 $t=0$ 处对 $t$ 求导, 并利用缩并是双线性运算这一事实, 即得所求结论.
\end{enumerate}
\end{description}
